% 诗歌类导言区
\usepackage{titlesec}
\usepackage{xeCJK}
\usepackage{fontspec,xunicode,xltxtra}
\usepackage{xpinyin}
\usepackage{tcolorbox}
\usepackage{geometry}
\usepackage{indentfirst}
\usepackage{pifont}
\usepackage[perpage,symbol*]{footmisc}
\usepackage{eso-pic}
\usepackage{graphicx} 
\usepackage{tikz}

\ExplSyntaxOn
\cs_set_protected:Npn \__xpinyin_make_pinyin_box:nnn #1#2#3
  {
    \__xpinyin_leavevmode:
    \hbox_overlap_right:n
      {
        \hbox_set:Nn \l__xpinyin_tmpa_box
          { \__xpinyin_CJKsymbol_hook: \__xpinyin_save_CJKsymbol:n {#2} }
        \hbox_set:Nn \l__xpinyin_tmpb_box
          {
            \l__xpinyin_pinyin_box_hook_tl
            \__xpinyin_select_font:
            \l__xpinyin_format_tl
            \clist_if_exist:cT { c__xpinyin_multiple_ #1 _clist }
              { \l__xpinyin_multiple_tl }
            {#3}
          }
        % 原第 164–171 行的 \dim_compare 和 \box_resize_to_wd_and_ht 已删除
        \box_move_up:nn { \l__xpinyin_vsep_tl }
          {
            \hbox_to_wd:nn { \box_wd:N \l__xpinyin_tmpa_box }
              { \tex_hss:D \box_use_drop:N \l__xpinyin_tmpb_box \tex_hss:D }
          }
      }
  }
\ExplSyntaxOff

% 设置背景图
\newcommand{\SetPageBackground}[1]{%
  \AddToShipoutPictureBG*{%
    \AtPageLowerLeft{%
    \begin{tikzpicture}[opacity=0.75]
        \node[inner sep=0pt, anchor=south west] at (0,0)
        {\includegraphics[width=\paperwidth,height=\paperheight,
                            keepaspectratio]{#1}};
    \end{tikzpicture}%
    }%
  }%
}

% 定义拼音使用的缺省字体
\setmainfont{Mona Sans Light}
% 定义 Century Schoolbook 字体
\newfontfamily\schoolbookfont{Century Schoolbook}
% 定义中文字体
\setCJKmainfont[BoldFont=STZhongsong]{KaiTi}
% 定义隶书字体族
\setCJKfamilyfont{beishu}{YouYuan} % 定义背诵框字体
\newcommand{\beishu}{\CJKfamily{beishu}} % 创建一个 \beishu 命令，方便调用
\xeCJKDeclareCharClass{CJK}{`0 -> `9}
\xeCJKsetup{AllowBreakBetweenPuncts=true}
\DefineFNsymbols{circled}{{\ding{192}}{\ding{193}}{\ding{194}}{\ding{195}}{\ding{196}}{\ding{197}}{\ding{198}}{\ding{199}}{\ding{200}}{\ding{201}}}
\setfnsymbol{circled}
\xpinyinsetup{ratio=0.48,hsep={.5em plus .1em minus .1em},vsep={1em}}
\setpinyin{一}{yi4}
\setpinyin{尽}{jin4}
\setpinyin{地}{di4}
\setpinyin{长}{chang2}
\setpinyin{朝}{zhao1}
\setpinyin{罗}{luo2}
\setpinyin{铺}{pu1}
\setpinyin{泊}{bo2}
\setpinyin{似}{si4}


% 定义背诵框环境
\newtcolorbox{beisongbox}{
    colback=yellow!10!white,   % 背景色
    colframe=orange!60!black,  % 边框色
    arc=2pt,                   % 圆角半径
    boxrule=0.6pt,             % 边框线粗细
    halign=center,             % 内容水平居中
    valign=center,             % 内容垂直居中
    nobeforeafter,             % 去掉框前后的额外间距
    boxsep=0pt,                % 内边距
    left=2pt,                  % 左内边距
    right=2pt,                 % 右内边距
    top=2pt,                   % 上内边距
    bottom=2pt,                % 下内边距
    fontupper=\beishu\zihao{3}\color{red!70!black}\ziju{-0.0}, % 框内字体、字号等
}
% 背诵框命令
\newcommand{\beisong}[1]{
    \begin{textblock*}{30mm}(15mm, 13mm)   % 宽 30mm，左上角在 (15mm, 13mm)  
    \begin{beisongbox}#1\end{beisongbox}
\end{textblock*}
}

% 章节标题字体格式
\titleformat{\chapter}{\centering\zihao{0}\bfseries}{}{0pt}{}
\titleformat{\section}{\zihao{-2}\bfseries}{ }{0pt}{}
\author{}
\date{}
